\documentclass[10pt,a4paper]{amsart}
\usepackage[margin=19mm]{geometry}
\usepackage{amsmath,amssymb,mathtools}
\usepackage[scr=rsfs,cal=euler]{mathalfa}
\usepackage{xfrac}
\usepackage[unicode,hidelinks,bookmarks=false]{hyperref}
\newcommand{\ie}{i.e.\ }
\newcommand{\cf}{cf.\ }
\newcommand{\resp}{respectively\ }
\newcommand{\Subsection}{\subsection}
\providecommand{\nobreakdash}{\nobreak\hbox{-}}
\setlength{\emergencystretch}{2em}
\multlinegap0pt
\usepackage{amssymb}
\usepackage{mathtools}
\usepackage{xcolor}
\newtheorem{lemma}{Lemma}
\newtheorem{proposition}{Proposition}
\newcommand{\BHint}[2]{\mathrm{BH}^{\mathrm{int}}_{#1,#2}}
\newcommand{\C}{\mathbb C}
\newcommand{\Indup}{\textcolor{black}{\operatorname{Ind}_{\uparrow}}}
\title{Subexponential Bohnenblust--Hille Inequalities\\on Finite Cyclic Groups}
\author{Daniel M. Pellegrino, Anselmo Raposo Jr}
\date{Source: arXiv:2608.05366v2 (CC BY 4.0)}
\begin{document}
\maketitle

\noindent\emph{Source and licence.} Subexponential Bohnenblust--Hille Inequalities\\on Finite Cyclic Groups. Version arXiv:2608.05366v2 (CC BY 4.0), distributed by arXiv under the Creative Commons Attribution 4.0 International licence, http://creativecommons.org/licenses/by/4.0/, as recorded in the arXiv licence metadata for this exact version and shown on the abstract page https://arxiv.org/abs/2608.05366v2. Full licence notice: \texttt{licenses/arxiv-2608.05366-CC-BY-4.0.txt}.\par\smallskip

\noindent\textit{Editorial scope.} Counted unit: the article's proposition prop:coarse-bound (source0.tex lines 1578--1589). The article's complete original proof of it is delivered with it (source0.tex lines 1591--1697, one \texttt{proof} environment of 107 lines). No mathematics is written, repaired or re-derived: the statement and the proof are the article's own lines, in the article's own notation, and no retained line is substituted or re-flowed.  Retained uncounted as the source-local context the proof needs: lines 748--754; lines 902--915; lines 1012--1014; lines 1040--1040; lines 1227--1230; lines 1232--1240; lines 1313--1334; lines 1356--1365; lines 1493--1498; lines 1558--1563.  Nothing was omitted from any retained span, so the package carries no omission record.  No retained line contains a citation, so the wrapper carries no bibliography and no bibliography entry.  No retained line contains a figure environment, an image-inclusion command or drawing code, so no image is delivered and the package declares no asset dependency.  Editorial formatting only, in addition: an AMS \texttt{amsart} preamble that restates the article's own declarations and macros used by the retained lines, namely the environments \texttt{lemma}, \texttt{proposition} on separate counters, together with the standard packages those lines need.  The retained environments print in this document as Lemma 1, Proposition 1 in order of appearance, whereas the article prints the same items with the numbering of its own full document.  The complete native source of the article as distributed in the arXiv e-print for this exact version is bundled unchanged as \texttt{sources/206919/source0.tex}.
\begin{equation}\label{eq:blei-body}
 \left(\sum_{\mathbf i\in I^d}|b_{\mathbf i}|^{p_d}\right)^{1/p_d}
 \leq
 \left(
 \prod_{\substack{S\subset[d]\\ |S|=m}}M_S(b)
 \right)^{1/\binom dm}.
\end{equation}

\begin{lemma}[Mixed polarization]\label{lem:qary-mixed-polarization}
For $0\leq m\leq d$ and $x,y\in C_q^N$,
\begin{equation}\label{eq:qary-mixed-polarization}
 |L_f(\zeta(x)^m,\zeta(y)^{d-m})|
 \leq A_{d,m}\|f\|_\infty,
\end{equation}
where
\begin{equation}\label{eq:Adm-sum-q}
 A_{d,m}
 :=\frac1{\binom dm}
 \sum_{r=0}^{\min\{m,d-m\}}
 4^r\binom d{2r}\binom{d-2r}{m-r}.
\end{equation}
\end{lemma}

\begin{equation}\label{eq:Adm-exact-q}
 A_{d,m}=\frac{\binom{2d}{2m}}{\binom dm}.
\end{equation}

\subsection{Orbit notation and normalization}\label{sec:orbits}

\begin{equation}\label{eq:atilde-C-bound}
 |\widetilde a_{\mathbf i_S\oplus\mathbf j_{S^c}}|
 \leq\binom dm|C^S_{\mathbf i_S,\mathbf j_{S^c}}|.
\end{equation}

\begin{lemma}[Mixed evaluation identity]\label{lem:mixed-evaluation}
Fix $S\subset[d]$ with $|S|=m$.  For every $x,y\in C_q^N$,
\begin{equation}\label{eq:mixed-evaluation}
 L_f(\zeta(x)^m,\zeta(y)^{d-m})
 =\sum_{\mathbf i_S\in\Indup(S)}
  \sum_{\mathbf j_{S^c}\in\Indup(S^c)}
 C^S_{\mathbf i_S,\mathbf j_{S^c}}x^{\mathbf i_S}y^{\mathbf j_{S^c}}.
\end{equation}
\end{lemma}

\begin{lemma}[Reduction to block estimates]\label{lem:blei-reduction-recurrence}
Let $2\leq m\leq d/2$, and let
$f\in\mathcal V_{\leq d}(C_q^N)$.  For $S\subset[d]$ with $|S|=m$, set
\begin{equation}\label{eq:QS}
 Q_S:=
 \left[
 \sum_{\mathbf i_S\in\Indup(S)}
 \left(
 \sum_{\mathbf j_{S^c}\in\Indup(S^c)}
 |\widetilde a_{\mathbf i_S\oplus\mathbf j_{S^c}}|^2
 \right)^{p_m/2}
 \right]^{1/p_m}.
\end{equation}
Then
\begin{equation}\label{eq:blei-canonical}
 \|\widehat f\|_{\ell_{p_d}}
 \leq
 \left(
 \prod_{\substack{S\subset[d]\\ |S|=m}}Q_S
 \right)^{1/\binom dm}.
\end{equation}
\end{lemma}

\begin{lemma}[Estimate of one block]\label{lem:one-block-estimate}
Under the hypotheses of Lemma~\ref{lem:blei-reduction-recurrence}, for every
$S\subset[d]$ with $|S|=m$,
\begin{equation}\label{eq:QS-final}
 Q_S\leq
 \binom dm
 \left(\frac{m+1}{m-1}\right)^{\gamma_q(d-m)}
 \BHint{m}{q}A_{d,m}\|f\|_\infty.
\end{equation}
\end{lemma}

\begin{lemma}[Interaction order one]\label{lem:interaction-one}
For every fixed $q\geq2$,
\begin{equation}\label{eq:B1-elementary}
 \BHint{1}{q}\leq 1+2\pi(q-1).
\end{equation}
\end{lemma}

\begin{lemma}[Endpoint $L_1$--$L_2$ estimate]\label{lem:endpoint-L1-L2}
If $g:C_q^N\to\C$ has interaction order at most $k$, then
\begin{equation}\label{eq:L1-L2-coarse}
 \|g\|_2\leq2^{3\gamma_q k}\|g\|_1.
\end{equation}
\end{lemma}

\begin{proposition}[Coarse exponential control]\label{prop:coarse-bound}
For every fixed $q\geq2$ and every integer $r\geq2$,
\begin{equation}\label{eq:coarse-exponential}
 \BHint{r}{q}
 \leq
 \BHint{1}{q}\,2^{3\gamma_q(r-1)}\binom{2r}{2}.
\end{equation}
Consequently,
\begin{equation}\label{eq:coarse-bound-log}
 \log\BHint{r}{q}=O_q(r).
\end{equation}
\end{proposition}

\begin{proof}
We use a singleton block directly; this is the endpoint version of the
coefficient decomposition used above, rather than an application of
Lemmas~\ref{lem:blei-reduction-recurrence}--\ref{lem:one-block-estimate}, whose
hypotheses require $m\geq2$.

Fix $S\subset[r]$ with $|S|=1$.  We use the notation of the orbit
construction in Section~\ref{sec:orbits} with $d$ replaced by $r$ and
$m=1$; in particular, $C^S$ denotes the corresponding two-block coefficient
array.  Blei's inequality
\eqref{eq:blei-body}, now with $m=1$ and hence $p_1=1$, gives
\begin{equation}\label{eq:coarse-Blei-singleton}
 \|\widehat f\|_{\ell_{p_r}}
 \leq
 \left(
 \prod_{|S|=1} Q_S^{(1)}
 \right)^{1/r},
\end{equation}
where
\[
 Q_S^{(1)}:=
 \sum_{\mathbf i_S\in\Indup(S)}
 \left(
 \sum_{\mathbf j_{S^c}\in\Indup(S^c)}
 |\widetilde a_{\mathbf i_S\oplus\mathbf j_{S^c}}|^2
 \right)^{1/2}.
\]
As in the proof of Lemma~\ref{lem:one-block-estimate}, the orbit-ratio bound
\eqref{eq:atilde-C-bound} gives
\begin{equation}\label{eq:coarse-orbit-singleton}
 Q_S^{(1)}
 \leq r
 \sum_{\mathbf i_S}
 \left(
 \sum_{\mathbf j_{S^c}}
 |C^S_{\mathbf i_S,\mathbf j_{S^c}}|^2
 \right)^{1/2}.
\end{equation}
For fixed $\mathbf i_S$, set
\[
 g_{\mathbf i_S}(y)
 :=\sum_{\mathbf j_{S^c}}
 C^S_{\mathbf i_S,\mathbf j_{S^c}}y^{\mathbf j_{S^c}}.
\]
The nonzero terms are distinct characters, so Parseval and
Lemma~\ref{lem:endpoint-L1-L2} yield
\[
 \left(
 \sum_{\mathbf j_{S^c}}
 |C^S_{\mathbf i_S,\mathbf j_{S^c}}|^2
 \right)^{1/2}
 =\|g_{\mathbf i_S}\|_2
 \leq
 2^{3\gamma_q(r-1)}\|g_{\mathbf i_S}\|_1.
\]
Summing in $\mathbf i_S$ and using Fubini on the normalized counting measure,
\begin{equation}\label{eq:coarse-Fubini-singleton}
 \sum_{\mathbf i_S}
 \left(
 \sum_{\mathbf j_{S^c}}
 |C^S_{\mathbf i_S,\mathbf j_{S^c}}|^2
 \right)^{1/2}
 \leq
 2^{3\gamma_q(r-1)}
 \sup_{y\in C_q^N}\sum_{\mathbf i_S}|g_{\mathbf i_S}(y)|.
\end{equation}
For fixed $y$, define
\[
 h_y(x):=L_f(\zeta(x),\zeta(y)^{r-1}).
\]
Lemma~\ref{lem:mixed-evaluation} gives
\[
 h_y(x)=\sum_{\mathbf i_S}g_{\mathbf i_S}(y)x^{\mathbf i_S}.
\]
The nonzero characters in this sum are distinct and have interaction order at
most one.  Since $p_1=1$, the definition of $\BHint{1}{q}$ therefore gives
\[
 \sum_{\mathbf i_S}|g_{\mathbf i_S}(y)|
 \leq \BHint{1}{q}\,\|h_y\|_\infty.
\]
Finally, Lemma~\ref{lem:qary-mixed-polarization} with $m=1$ yields
\[
 \sup_y\|h_y\|_\infty\leq A_{r,1}\|f\|_\infty.
\]
Combining this with \eqref{eq:coarse-orbit-singleton} and
\eqref{eq:coarse-Fubini-singleton}, we obtain, for every singleton $S$,
\[
 Q_S^{(1)}
 \leq
 r\,2^{3\gamma_q(r-1)}\BHint{1}{q}A_{r,1}\|f\|_\infty.
\]
Insert this bound into \eqref{eq:coarse-Blei-singleton}.  Since the right-hand
side is independent of $S$,
\[
 \|\widehat f\|_{\ell_{p_r}}
 \leq
 r\,2^{3\gamma_q(r-1)}\BHint{1}{q}A_{r,1}\|f\|_\infty.
\]
Taking the supremum over $N$ and nonzero
$f\in\mathcal V_{\leq r}(C_q^N)$, and using
\[
 rA_{r,1}=\binom{2r}{2}
\]
from \eqref{eq:Adm-exact-q}, proves \eqref{eq:coarse-exponential}.  Taking
logarithms and using Lemma~\ref{lem:interaction-one} gives
\eqref{eq:coarse-bound-log}.
\end{proof}

\end{document}
